\documentclass[11pt,a4paper]{article}

% ----------------------------------------------------------------------
% Packages
% ----------------------------------------------------------------------
\usepackage[margin=1in]{geometry}
\usepackage[T1]{fontenc}
\usepackage[utf8]{inputenc}
\usepackage{lmodern}
\usepackage{amsmath,amssymb}
\usepackage{booktabs}
\usepackage{longtable}
\usepackage{array}
\usepackage{tabularx}
\usepackage{makecell}
\usepackage{enumitem}
\usepackage{xcolor}
\usepackage{hyperref}
\usepackage{titlesec}
\usepackage{fancyhdr}
\usepackage{tcolorbox}
\usepackage{graphicx}
\usepackage{float}
\usepackage{caption}
\usepackage{tikz}
\usetikzlibrary{arrows.meta, positioning, shapes.geometric, calc}

% ----------------------------------------------------------------------
% Hyperlinks and colors
% ----------------------------------------------------------------------
\hypersetup{
    colorlinks=true,
    linkcolor=blue!60!black,
    urlcolor=blue!60!black,
    citecolor=blue!60!black
}

% ----------------------------------------------------------------------
% Custom commands
% ----------------------------------------------------------------------
\newcommand{\code}[1]{\texttt{#1}}
\newcommand{\GF}[1]{\mathrm{GF}(#1)}
\newcommand{\reporttitle}{SemiCON-FabricATE}
\newcommand{\reportsubtitle}{AXI4-to-PCIe 6 Link RTL: Technical Report}

% Simple highlighted box for key notes
\newtcolorbox{keybox}{
    colback=blue!5,
    colframe=blue!50!black,
    boxrule=0.5pt,
    arc=2pt,
    left=6pt,
    right=6pt,
    top=6pt,
    bottom=6pt
}

% ----------------------------------------------------------------------
% Header / footer
% ----------------------------------------------------------------------
\pagestyle{fancy}
\fancyhf{}
\fancyhead[L]{\reporttitle}
\fancyhead[R]{SARCathon 2026}
\fancyfoot[C]{\thepage}
\renewcommand{\headrulewidth}{0.4pt}

% ----------------------------------------------------------------------
% Document
% ----------------------------------------------------------------------
\begin{document}

% ======================================================================
% Title Page
% ======================================================================
\begin{titlepage}
    \centering
    \vspace*{1.2cm}

    {\scshape\Large SARCathon 2026 -- Semiconductor Challenge\par}
    \vspace{0.4cm}
    {\large Design the Bridge. Connect the Future.\par}

    \vspace{2.2cm}

    {\Huge\bfseries \reporttitle\par}
    \vspace{0.6cm}
    {\Large \reportsubtitle\par}

    \vspace{1.5cm}

    \begin{tabular}{>{\bfseries}ll}
        Project: & SemiCON-FabricATE \\
        Challenge: & SARCathon 2026 SemiCON Hackathon \\
        Repository: & \code{github.com/cyberaionics/SemiCON-FabricATE} \\
        Report date: & 2026-09-11 \\
        Assessment basis: & Repository RTL, testbenches, scripts, \\
        & recorded logs, design documentation, and \\
        & the supplied problem-statement PDF \\
    \end{tabular}

    \vfill

    {\large Technical Report\par}
    \vspace{0.3cm}
    {\large Stage 1 Qualifier \par}
\end{titlepage}

% ======================================================================
% Table of Contents
% ======================================================================
\tableofcontents
\newpage

% ======================================================================
% 1. Executive Summary
% ======================================================================
\section{Executive Summary}

The supplied challenge asks teams to develop an AXI4-to-PCIe interface, verify its RTL,
validate it against meaningful workloads, synthesize it, and ultimately carry the design
toward physical implementation and GDSII. The repository delivers a substantial
transmit-side RTL prototype in two stages:

\begin{enumerate}[leftmargin=*]
    \item \textbf{Stage 1:} two 512-bit AXI4-Stream inputs are buffered, arbitrated without
    packet interleaving, compacted into application-data frames, protected with an
    eight-byte CRC and six-byte FEC field, and emitted as eight 256-bit transfers per frame.
    \item \textbf{Stage 2:} an independent 256-bit protocol-to-PHY digital path stripes data
    across four lanes, scrambles each lane, applies Gray/PAM4 symbol mapping, and reverses
    the operation at the receiver under control of a link-training state machine.
\end{enumerate}

The implementation is synthesizable SystemVerilog and has strong simulation evidence.
The current regression executed successfully with Icarus Verilog 13.0. It covered CRC and
FEC unit behavior, error-injection models, arbitration, interconnect traffic, packet
framing, output backpressure, packet splitting, sparse \code{TKEEP}, resets, and Stage 2
PHY/LTSSM behavior.

\begin{keybox}
The design should be described accurately as a \textbf{verified transmit/framing and
digital-PHY prototype}, not as a complete PCIe 6.0 AXI4 bridge. It does not implement real
AXI memory-mapped transactions, PCIe TLP/DLP encoding, receive-side CRC/FEC hardware,
replay or credit management, PCIe ordered sets, analog SerDes behavior, application
workload validation, FPGA fitting, timing closure, OpenROAD/OpenLane implementation, or
GDSII generation. Stage 2 is currently not wired to the Stage 1 frame output.
\end{keybox}

% ======================================================================
% 2. Project Identity and Requirement Interpretation
% ======================================================================
\section{Project Identity and Requirement Interpretation}

\subsection{Project Name}
\textbf{SemiCON-FabricATE} --- AXI4-to-PCIe 6 Link RTL Prototype.

\subsection{Challenge Context}
The challenge is SARCathon 2026 SemiCON Hackathon, with the theme
``Design the Bridge. Connect the Future.'' The supplied problem statement defines a
two-stage path:

\begin{itemize}[leftmargin=*]
    \item \textbf{Stage 1 -- Qualifier:} understand the architecture, develop synthesizable
    RTL, build functional testbenches, verify reset/error behavior, and provide simulation
    results with a short design report.
    \item \textbf{Stage 2 -- Grand Finale:} complete RTL integration, exhaustive functional
    verification, application-level workload validation, RTL optimization, synthesis,
    physical design through OpenROAD/OpenLane, and final GDSII.
\end{itemize}

\subsection{Report Basis}
This report is based on the repository RTL, testbenches, scripts, recorded logs, design
documentation, and the supplied problem-statement PDF
\code{6a9861f9d0ed3\_semicon\_ps.pdf}. It also uses the supplied markdown technical report
as the primary content source and the \emph{Protocol Playground Proposal} as a formatting
reference for title page, section numbering, tables, and project-style presentation.

\subsection{Requirement Interpretation}
The problem statement asks for an AXI4-to-PCIe 6 link based on the reference architecture:
AXI4/AXI4-Stream interface, interconnect layer, link layer with packetization and
CRC/FEC concepts, and PCIe protocol-to-PHY layer with lane striping, Gray coding, and
PAM4 mapping.

\medskip
\noindent
\begin{tabularx}{\textwidth}{>{\raggedright\arraybackslash}p{0.30\textwidth} X}
\toprule
\textbf{Challenge Area} & \textbf{Interpretation for This Report} \\
\midrule
AXI4 / AXI4-Stream interface & Implemented as 512-bit AXI4-Stream-like ingress with
ready/valid handshakes and source identification. No memory-mapped AXI read/write channel
is present. \\
Interconnect layer & Implemented as a two-to-one packet-held round-robin arbiter with
per-source FIFOs and a two-beat output FIFO. \\
Link layer & Implemented as a byte-compacting frame packetizer, eight-byte CRC generator,
and six-byte FEC generator producing fixed 256-byte frames. \\
PCIe protocol-to-PHY layer & Implemented as an independent 256-bit digital PHY prototype
with four-lane striping, 23-bit lane scramblers, Gray coding, and two-bit PAM4 symbol
mapping. \\
Synthesis and physical design & Generic Yosys structural checks are present. No
technology-specific synthesis, OpenROAD/OpenLane flow, timing closure, or GDSII exists. \\
Workload validation & No network-packet or ML-inference workload generators/results are
present. \\
\bottomrule
\end{tabularx}

% ======================================================================
% 3. Delivered Scope and Top-Level Data Flow
% ======================================================================
\section{Delivered Scope and Top-Level Data Flow}

\subsection{Stage 1 Tested Path}
The tested Stage 1 path is:

\begin{verbatim}
AXI-Stream source 0 (512 bits)
        \
         -> input FIFO 0 ----\
                              packet-held round-robin arbiter
         -> input FIFO 1 ----/          |
                                         v
                                  two-beat output FIFO
                                         |
                                         v
                              byte collection and framing
                                         |
                                         v
                                    CRC generation
                                         |
                                         v
                                    FEC generation
                                         |
                                         v
                              eight 256-bit output transfers
\end{verbatim}

The principal integrated top level is \code{rtl/axis\_link\_tx.sv}. Its input boundary is
two 512-bit AXI4-Stream-like interfaces. Its output boundary is a 256-bit ready/valid
stream with frame metadata.

\subsection{Stage 2 Independent Path}
The independent Stage 2 path is:

\begin{verbatim}
256-bit input
  -> byte striping over four lanes
  -> independent 23-bit lane scramblers
  -> Gray coding
  -> two-bit digital PAM4 symbols
  -> inverse RX mapping, descrambling, and destriping
  -> 256-bit output
\end{verbatim}

The Stage 2 top level is \code{rtl/protocol\_to\_phy\_top.sv}. Its LTSSM controls whether
the TX and RX datapaths are enabled.

\subsection{Integration Status}
No direct RTL connection currently exists between \code{axis\_link\_tx} and
\code{protocol\_to\_phy\_top}. Integration remains a future milestone.

% ======================================================================
% 4. Stage 1 Architecture
% ======================================================================
\section{Stage 1 Architecture}

\subsection{AXI4-Stream Ingress and Arbitration}
\code{rtl/axis\_interconnect\_2to1.sv} instantiates one FIFO per source and a two-beat
output FIFO. The FIFO preserves complete input beats and their \code{TKEEP}, \code{TLAST},
and source identity. It does not compact bytes; compaction occurs later in the packetizer.

\code{rtl/packet\_rr\_arbiter.sv} implements packet-held round-robin arbitration:

\begin{itemize}[leftmargin=*]
    \item ownership is acquired when a source presents a selectable beat;
    \item ownership remains fixed during stalls and source gaps;
    \item ownership is released only after the selected \code{TLAST} beat is accepted;
    \item the next preference alternates after a completed packet; and
    \item a missing preferred source does not prevent the other source from progressing.
\end{itemize}

This prevents packet interleaving, but fairness is packet-based and assumes packets
eventually terminate. A source that never asserts \code{TLAST} can block the other source
indefinitely because there is no timeout or abort protocol.

The stream convention is byte lane zero at \code{TDATA[7:0]}, with \code{TKEEP[0]}
qualifying that byte. Invalid bytes are retained through the interconnect and discarded
only by the packetizer. FIFO readiness depends on current occupancy, so a full FIFO does
not recover in the same cycle as a pop; this is functionally conservative but reduces
throughput, especially at depth one.

\subsection{Frame Packetizer}
\code{rtl/axis\_frame\_packetizer.sv} uses the following phases:

\begin{table}[H]
\centering
\caption{Packetizer phases}
\begin{tabular}{>{\bfseries}ll}
\toprule
Phase & Function \\
\midrule
\code{COLLECT} & Accept and store one 512-bit input beat. \\
\code{SCAN} & Examine all 64 byte lanes, one lane per cycle, and append only qualified bytes. \\
\code{FINISH} & Determine whether the packet ended, whether the frame is full, and whether another frame is required. \\
\code{CRC} & Process the first 242 frame bytes. \\
\code{FEC} & Process the first 250 frame bytes, including the generated CRC. \\
\code{EMIT} & Present eight 256-bit output transfers with stable metadata. \\
\bottomrule
\end{tabular}
\end{table}

The packetizer stores one input beat and a 236-byte application payload region. It does not
combine application packets in one frame. Long packets are divided into multiple frames. An
empty packet produces one zero-byte frame. Exact 236-byte packets do not produce an
unwanted extra empty frame, even if a later all-zero-\code{TKEEP} beat carries \code{TLAST}.

During output backpressure, \code{TDATA}, \code{TKEEP}, \code{TLAST}, source, byte count,
and packet flags remain stable. Reset aborts any partial collection, parity calculation, or
frame transmission.

\subsection{Frame Format}
Every emitted frame is 256 bytes and is transferred as eight 32-byte beats.

\begin{table}[H]
\centering
\caption{Frame format}
\begin{tabular}{>{\raggedleft\arraybackslash}p{0.18\textwidth} >{\raggedleft\arraybackslash}p{0.12\textwidth} p{0.38\textwidth} p{0.20\textwidth}}
\toprule
\textbf{Byte offsets} & \textbf{Size} & \textbf{Contents} & \textbf{Protection} \\
\midrule
0--235 & 236 bytes & Compacted application bytes followed by zero padding & CRC and FEC \\
236--241 & 6 bytes & Reserved data-link region, currently zero & CRC and FEC \\
242--249 & 8 bytes & CRC0 through CRC7 & FEC only \\
250--255 & 6 bytes & \code{C1, C2, C0, P1, P2, P0} & Transmitted FEC \\
\bottomrule
\end{tabular}
\end{table}

Output \code{TKEEP} is all ones for every valid transfer. Padding, reserved bytes, CRC, and
FEC are transmitted bytes, so the output is a fixed-size frame rather than a variable-size
final transfer.

Frame metadata is repeated on all eight transfers:

\begin{table}[H]
\centering
\caption{Frame metadata sidebands}
\begin{tabular}{>{\raggedleft\arraybackslash}p{0.22\textwidth} >{\raggedleft\arraybackslash}p{0.08\textwidth} p{0.58\textwidth}}
\toprule
\textbf{Signal} & \textbf{Width} & \textbf{Meaning} \\
\midrule
\code{m\_axis\_source} & 1 & Original source number, 0 or 1 \\
\code{m\_frame\_bytes} & 8 & Application bytes represented by this frame, 0--236 \\
\code{m\_packet\_start} & 1 & First frame of an application packet \\
\code{m\_packet\_end} & 1 & Final frame of an application packet \\
\code{m\_axis\_tlast} & 1 & Final transfer of the 256-byte frame, independent of packet end \\
\bottomrule
\end{tabular}
\end{table}

These sidebands are internal prototype metadata. They are not inserted into the protected
bytes and would need a defined downstream representation in a complete link.

% ======================================================================
% 5. CRC and FEC Implementation
% ======================================================================
\section{CRC and FEC Implementation}

\subsection{Eight-Byte CRC}
\code{rtl/flit\_crc8.sv} computes an eight-byte Reed--Solomon CRC over frame bytes 0--241:

\begin{itemize}[leftmargin=*]
    \item field: \(\GF{256}\);
    \item primitive polynomial: \code{0x12B}, or \(x^8+x^5+x^3+x+1\);
    \item field alpha: \code{0x02};
    \item generator coefficients: \code{01 D5 68 FE D5 33 41 4D 69};
    \item initial state: zero;
    \item reflection: none;
    \item final complement: none; and
    \item byte order: \code{crc[7:0]} is CRC0 and \code{crc[63:56]} is CRC7.
\end{itemize}

The CRC accumulator accepts one enabled byte per cycle. The caller controls clearing, byte
count, and frame boundaries. The packetizer performs the complete 242-cycle CRC pass before
starting FEC.

The Python reference in \code{sim/crc\_vectors.py} derives the generator through polynomial
operations rather than copying the RTL recurrence. This provides an independent oracle for
known vectors and integrated frames.

\subsection{Six-Byte FEC}
\code{rtl/flit\_fec6.sv} computes six ECC bytes over frame bytes 0--249, including the
completed CRC:

\begin{itemize}[leftmargin=*]
    \item field: \(\GF{256}\);
    \item primitive polynomial: \code{0x11D}, or \(x^8+x^4+x^3+x^2+1\);
    \item three interleaved groups selected by byte index modulo three;
    \item group 0 has 84 transmitted symbols;
    \item groups 1 and 2 have 83 transmitted symbols plus an implicit zero; and
    \item each group generates a weighted check value \(C\) and an XOR parity value \(P\).
\end{itemize}

The transmitted mapping is:

\begin{table}[H]
\centering
\caption{FEC transmitted mapping}
\begin{tabular}{>{\raggedleft\arraybackslash}p{0.15\textwidth} l}
\toprule
\textbf{Frame byte} & \textbf{Value} \\
\midrule
250 & \code{C1} \\
251 & \code{C2} \\
252 & \code{C0} \\
253 & \code{P1} \\
254 & \code{P2} \\
255 & \code{P0} \\
\bottomrule
\end{tabular}
\end{table}

The FEC pass consumes 250 enabled cycles after CRC insertion. The RTL is an encoder only.
The correction model in \code{sim/fec\_vectors.py} is verification software and is not a
synthesized receive datapath.

% ======================================================================
% 6. Stage 2 Digital PHY Prototype
% ======================================================================
\section{Stage 2 Digital PHY Prototype}

\subsection{Datapath}
The Stage 2 modules implement a reversible digital transformation:

\begin{enumerate}[leftmargin=*]
    \item \code{lane\_striper.sv} distributes bytes over four lanes.
    \item \code{lane\_scrambler.sv} applies an independent 23-bit LFSR-based operation on each lane.
    \item \code{gray\_pam4\_tx.sv} maps pairs of bits through Gray coding to two-bit PAM4 symbols.
    \item \code{gray\_pam4\_rx.sv} reverses the symbol mapping and Gray coding.
    \item \code{lane\_destriper.sv} reconstructs the original 256-bit word.
\end{enumerate}

\code{protocol\_to\_phy\_tx.sv} and \code{protocol\_to\_phy\_rx.sv} provide ready/valid
buffering around the transform. RX output is held until accepted, and the TX path is gated
by link state.

\subsection{LTSSM}
\code{rtl/ltssm.sv} models Detect, Polling, Configuration, L0, Recovery, Disabled, and Hot
Reset behavior. It exposes \code{link\_up}, TX/RX enable controls, training/recovery
indicators, the current state, and a state-change pulse.

The Stage 2 testbench exercises invalid link width blocking, link bring-up, recovery
requests, PHY-error recovery, disabled/retraining, and hot reset. One implementation detail
should be addressed before claiming a parameter-complete LTSSM:
\code{DETECT\_WAIT\_CYCLES} is exposed as a parameter but the current Detect path does not
use it to delay the next state.

% ======================================================================
% 7. Verification Strategy and Results
% ======================================================================
\section{Verification Strategy and Results}

\subsection{Current Executable Validation}
On 2026-09-13, the repository regression was run with Python 3 and Icarus Verilog 13.0
using:

\begin{verbatim}
python3 sim/run_all.py
\end{verbatim}

The run completed successfully. The key results were:

\begin{table}[H]
\centering
\caption{Stage 1 regression results}
\begin{tabular}{>{\raggedright\arraybackslash}p{0.28\textwidth} p{0.62\textwidth}}
\toprule
\textbf{Area} & \textbf{Result} \\
\midrule
FEC RTL & 2,032 vectors, including 2,000 basis-bit vectors, enable gaps, clear priority, and three reset positions: PASS \\
FEC software correction & 2,302 injected-error frames restored: 2,048 bit flips and 254 three-byte bursts: PASS \\
CRC RTL & 32 vectors, enable gaps, clear/reset, and 2,128 corrupted codewords rejected: PASS \\
Arbiter & Stalled first/final beat, source gaps, round robin, and reset: PASS \\
Interconnect & 10 randomized configurations/runs, 10,800 delivered beats overall: PASS \\
Stage 1 packetizer/link & Standalone and integrated runs across depths 1, 3, and 4, 840 packets and 2,716 frames in the recorded regression set: PASS \\
Waveform run & Integrated \code{sim/link\_tx.vcd} generated successfully: PASS \\
\bottomrule
\end{tabular}
\end{table}

The integrated checker covers:

\begin{itemize}[leftmargin=*]
    \item empty packets;
    \item sparse, null, and partial \code{TKEEP} masks;
    \item packet lengths around 236-byte boundaries;
    \item long packets split across frames;
    \item source ownership across packet boundaries;
    \item backpressure at all eight output positions;
    \item output stability while stalled;
    \item reset during collection, CRC, FEC, and output; and
    \item fresh traffic after reset.
\end{itemize}

The measured architecture is intentionally low-throughput. A 512-bit input beat is scanned
one byte lane per cycle, so collection alone introduces approximately 66 cycles between
input handshakes. Each frame adds 242 CRC cycles, 250 FEC cycles, and at least eight output
handshake cycles, with stalls extending the latency.

\subsection{Stage 2 Executable Validation}
The three Stage 2 benches were also compiled and run directly with Icarus 13.0. The
integrated bench completed with 113 checks/transactions and reported:

\begin{itemize}[leftmargin=*]
    \item invalid link width correctly blocked L0;
    \item LTSSM reached L0;
    \item directed TX-to-four-lane-PAM4-to-RX loopback passed;
    \item RX ready/valid backpressure passed;
    \item injected symbol corruption was observable at RX;
    \item explicit Recovery and \code{phy\_error} both returned to L0;
    \item 100 randomized end-to-end transactions passed;
    \item Disabled returned toward Detect;
    \item retraining transferred data; and
    \item Hot Reset returned toward Detect.
\end{itemize}

The block-level bench passed striping/destriping, Gray/PAM4 inverse mapping, scrambler round
trips, and a second directed pattern. The standalone LTSSM bench passed L0,
Recovery-to-L0, and Disabled-to-Detect checks.

Icarus emitted expected support warnings for constant selects in \code{always\_*} blocks in
the scrambler and missing explicit time units in several modules. These did not cause a test
failure, but adding explicit time units and considering \code{always\_comb}/\code{always\_ff}
portability would make tool behavior clearer.

\subsection{Verification Limits}
The tests are strong directed and randomized simulation evidence, but they are not
exhaustive proof. In particular:

\begin{itemize}[leftmargin=*]
    \item no formal properties or equivalence checking are included;
    \item receive-side CRC checking and FEC correction are not RTL-tested because they do not exist as hardware;
    \item no real PCIe transaction workload is generated;
    \item the corruption tests do not establish complete multi-error detection guarantees;
    \item Stage 1 and Stage 2 are not verified as one connected end-to-end design; and
    \item the Stage 2 result logs referenced by \code{docs/VERIFICATION.md} are absent from this checkout even though the tests can be run directly.
\end{itemize}

% ======================================================================
% 8. Synthesis and Implementation Status
% ======================================================================
\section{Synthesis and Implementation Status}

The checked-in scripts provide generic Yosys structural checks:

\begin{itemize}[leftmargin=*]
    \item \code{synthesis/check\_link.ys} targets the Stage 1 \code{axis\_link\_tx} hierarchy;
    \item \code{synthesis/run\_yosys.ys} targets the Stage 2 \code{protocol\_to\_phy\_top} hierarchy; and
    \item \code{check -assert} is included after synthesis.
\end{itemize}

Recorded repository evidence reports:

\begin{table}[H]
\centering
\caption{Recorded generic synthesis results}
\begin{tabular}{>{\raggedright\arraybackslash}p{0.34\textwidth} p{0.56\textwidth}}
\toprule
\textbf{Design} & \textbf{Recorded result} \\
\midrule
Stage 1 CRC/FEC-integrated path & Yosys structural synthesis passed; 38,148 unmapped generic primitive cells; no inferred latches; \code{check -assert} passed \\
Stage 2 digital PHY path & Yosys structural synthesis reported 7,758 primitive cells and 4,171 wires; no inferred latches; \code{check -assert} passed \\
\bottomrule
\end{tabular}
\end{table}

These numbers are generic structural counts, not standard-cell area, FPGA utilization,
timing, power, or a technology-specific result. The Stage 1 FIFOs and Stage 2 small
scrambler memories are expected to map to registers in this implementation.

The current machine does not have \code{yosys} on \code{PATH}, so the recorded synthesis
results could not be independently regenerated during this review. Additionally,
\code{synthesis/run\_yosys.sh} invokes a nonexistent \code{synthesis/run\_yosys.py}; the
checked-in \code{.ys} script must be invoked directly or the wrapper must be repaired.

No OpenROAD/OpenLane configuration, floorplan, placement, routing, parasitic extraction,
timing report, signoff report, or GDSII was found. No FPGA fit or hardware test evidence was
found.

% ======================================================================
% 9. Traceability to the Challenge
% ======================================================================
\section{Traceability to the Challenge}

\begin{longtable}{>{\raggedright\arraybackslash}p{0.34\textwidth} >{\raggedright\arraybackslash}p{0.42\textwidth} >{\raggedright\arraybackslash}p{0.18\textwidth}}
\caption{Traceability to SARCathon 2026 requirements}\\
\toprule
\textbf{Challenge requirement} & \textbf{Repository evidence} & \textbf{Status} \\
\midrule
\endfirsthead
\toprule
\textbf{Challenge requirement} & \textbf{Repository evidence} & \textbf{Status} \\
\midrule
\endhead
Develop synthesizable RTL & \code{rtl/*.sv}, Stage 1 and Stage 2 tops & Addressed for prototype scope \\
Provide functional testbench & \code{tb/*.sv}, independent Python reference models & Addressed \\
Verify reset behavior & Unit, interconnect, packetizer, CRC/FEC, and LTSSM tests & Addressed for implemented blocks \\
Verify error behavior & CRC corruption rejection, FEC correction model, PHY symbol corruption visibility & Partially addressed; no receive FEC/CRC hardware \\
Verify key AXI/PCIe read/write transactions & AXI4-Stream beat movement is tested; no memory-mapped AXI or PCIe TLP transactions & Not addressed \\
Demonstrate simulation results & Current Icarus regression and recorded logs & Addressed \\
Optimize for performance/area/reliability & Functional baseline and generic structural synthesis only & Partially addressed \\
Test network packet workloads & No workload model or results & Not addressed \\
Test ML-inference workloads & No workload model or results & Not addressed \\
Synthesize and analyze implementation & Generic Yosys scripts and recorded structural counts & Partially addressed \\
Complete OpenROAD/OpenLane flow & No flow configuration or reports & Not addressed \\
Generate GDSII & No GDSII artifact & Not addressed \\
FPGA-tested bonus & No FPGA project, fit, or board evidence & Not addressed \\
\bottomrule
\end{longtable}

% ======================================================================
% 10. Limitations and Engineering Risks
% ======================================================================
\section{Limitations and Engineering Risks}

\begin{longtable}{>{\raggedright\arraybackslash}p{0.08\textwidth} >{\raggedright\arraybackslash}p{0.44\textwidth} >{\raggedright\arraybackslash}p{0.38\textwidth}}
\caption{Engineering risks and mitigation considerations}\\
\toprule
\textbf{No.} & \textbf{Risk / limitation} & \textbf{Impact / mitigation} \\
\midrule
\endfirsthead
\toprule
\textbf{No.} & \textbf{Risk / limitation} & \textbf{Impact / mitigation} \\
\midrule
\endhead
1 & Not a complete PCIe bridge. The frame payload is raw application data. TLP headers, addresses, attributes, tags, DLP semantics, sequence numbers, replay, credits, flow control, and completion handling are absent. & Treat as transmit-side prototype; add protocol boundary before claiming PCIe compliance. \\
2 & Transmit-only protection. CRC and FEC encoders exist, but receive-side checking and correction are not synthesized RTL. & Implement RX CRC/FEC check and correction. \\
3 & Stage 2 is disconnected. The PHY top level begins at an independent 256-bit interface rather than consuming \code{axis\_link\_tx} frames. & Add adapter and end-to-end testbench. \\
4 & No read path. The design does not provide host-to-accelerator or PCIe-to-AXI return traffic. & Add request/completion descriptors and return traffic. \\
5 & Packet termination dependency. A missing \code{TLAST} can block the other source indefinitely. & Add timeout/abort or protocol-level termination guarantee. \\
6 & Low throughput. One-byte-per-cycle scanning and sequential CRC/FEC passes create large per-frame latency. & Compact multiple bytes per cycle; pipeline CRC/FEC. \\
7 & Register-heavy buffering. Shallow FIFO memories and the 236-byte frame store may consume substantial resources. & Use inferred or instantiated target memories where appropriate. \\
8 & Parameter checks are implicit. Invalid widths and illegal FIFO depths are documented but not guarded by explicit elaboration-time assertions. & Add elaboration-time parameter assertions. \\
9 & LTSSM parameter gap. \code{DETECT\_WAIT\_CYCLES} is declared but not used in the Detect transition behavior. & Use parameter in Detect path or remove it. \\
10 & Digital PHY abstraction. PAM4 symbols are represented as digital two-bit values; physical SerDes and compliance behavior are outside the model. & State abstraction clearly; add analog/SerDes model if required. \\
11 & Tooling inconsistency. The Linux Yosys wrapper references a missing Python script, and documented Stage 2 result artifacts are absent. & Repair wrapper and archive reproducible logs. \\
12 & No physical-design evidence. Generic synthesis is not a substitute for timing closure, power analysis, placement/routing, or GDSII signoff. & Complete OpenROAD/OpenLane flow. \\
\bottomrule
\end{longtable}

% ======================================================================
% 11. Recommended Next Steps
% ======================================================================
\section{Recommended Next Steps}

\subsection{Priority 1: Define and Integrate the Protocol Boundary}
\begin{enumerate}[leftmargin=*]
    \item Specify a real transaction model: AXI4 memory-mapped or a clearly constrained AXI4-Stream transaction layer.
    \item Add request/completion descriptors, address fields, tags, packet attributes, and return traffic.
    \item Define how frame metadata is encoded or transported downstream.
    \item Connect \code{axis\_link\_tx} to the Stage 2 256-bit input with a documented adapter and end-to-end testbench.
\end{enumerate}

\subsection{Priority 2: Complete Receive and Link-Layer Protection}
\begin{enumerate}[leftmargin=*]
    \item Implement RTL CRC checking and FEC syndrome/correction logic on the receive path.
    \item Verify correction followed by CRC validation for one-error-per-group and multi-error cases.
    \item Add sequence/replay handling, credit/flow control, and error reporting appropriate to the target protocol.
\end{enumerate}

\subsection{Priority 3: Improve Throughput and Resource Use}
\begin{enumerate}[leftmargin=*]
    \item Compact multiple input bytes per cycle instead of scanning one lane per cycle.
    \item Pipeline or parallelize CRC/FEC processing where the target clock and area budget justify it.
    \item Replace register-heavy FIFOs with inferred or instantiated target memories where appropriate.
    \item Add explicit parameter assertions and measure latency, throughput, area, and power under representative traffic.
\end{enumerate}

\subsection{Priority 4: Make Verification Challenge-Complete}
\begin{enumerate}[leftmargin=*]
    \item Add constrained-random transaction traffic for network packet and ML tensor workloads.
    \item Add assertions for ready/valid stability, packet ownership, frame boundaries, reset recovery, and protocol invariants.
    \item Add formal checks for FIFO safety, no packet interleaving, and output stability under arbitrary backpressure.
    \item Store reproducible Stage 2 logs and waveforms in the expected locations or update the documentation to match the artifact policy.
\end{enumerate}

\subsection{Priority 5: Complete Implementation Flow}
\begin{enumerate}[leftmargin=*]
    \item Repair \code{synthesis/run\_yosys.sh} or replace it with a direct Yosys wrapper.
    \item Select a target FPGA or open PDK and establish constraints.
    \item Add OpenROAD/OpenLane configuration, timing constraints, floorplanning, and signoff checks.
    \item Generate and archive reports for area, timing, power, DRC/LVS, and GDSII.
\end{enumerate}

% ======================================================================
% 12. Reproduction Commands
% ======================================================================
\section{Reproduction Commands}

\subsection{Stage 1 Regression}
\begin{verbatim}
python3 sim/run_all.py
\end{verbatim}

or:

\begin{verbatim}
./sim/run_all_crc.sh
\end{verbatim}

\subsection{Stage 2 Integrated Simulation}
\begin{verbatim}
mkdir -p sim/build
iverilog -g2012 -Wall -s tb_integrated -o sim/build/phy_integrated_sim \
  rtl/phy_pkg.sv rtl/ltssm.sv rtl/lane_striper.sv rtl/lane_destriper.sv \
  rtl/lane_scrambler.sv rtl/gray_pam4_tx.sv rtl/gray_pam4_rx.sv \
  rtl/protocol_to_phy_tx.sv rtl/protocol_to_phy_rx.sv \
  rtl/protocol_to_phy_top.sv tb/tb_integrated.sv
vvp sim/build/phy_integrated_sim
\end{verbatim}

The block-level and standalone LTSSM benches use the same source list pattern and are listed
in \code{sim/run\_all.ps1}.

\subsection{Generic Synthesis, When Yosys Is Installed}
\begin{verbatim}
yosys -l synthesis/link_structural_results.log synthesis/check_link.ys
yosys -l synthesis/phy_structural_results.log synthesis/run_yosys.ys
\end{verbatim}

% ======================================================================
% 13. Conclusion
% ======================================================================
\section{Conclusion}

SemiCON-FabricATE provides a coherent, synthesizable, and well-tested transmit-side
baseline. Its strongest contribution is the complete simulated Stage 1 data path:
packet-preserving arbitration, byte-compacting frame assembly, CRC/FEC generation, metadata
propagation, and robust ready/valid behavior under stalls and resets. The independent
Stage 2 prototype extends the work into a digital four-lane PAM4-style path with link-state
control and meaningful randomized verification.

The evidence supports claiming a functional RTL prototype and a strong Stage 1 simulation
result. It does not support claiming a complete AXI4-to-PCIe bridge, PCIe compliance, full
workload validation, FPGA implementation, physical design completion, or GDSII. The most
important engineering step is to close the boundary between the tested prototypes and then
add the missing transaction, receive, implementation, and workload-validation layers
required by the challenge statement.

\end{document}